\chapter{Related Work}
\label{chap:related}

\section{From Price-Only to Text-Aware Forecasting}

Research on stock prediction did not begin with sentiment. Early systems relied on prices, returns, volume, and technical indicators because these variables are directly observable and easy to align temporally. The main weakness of that tradition is explanatory rather than computational. Numerical series encode market state, but they do not encode the semantic content of the events that moved the market. Once researchers began to treat news and social media as external information channels, forecasting became a multimodal problem rather than a purely numerical one. \cite{fischer2018lstm, nelson2017lstm, sezer2020dlreview}

The shift from price-only to text-aware prediction followed a recognizable sequence. Schumaker and Chen's AZFinText system used breaking financial news to support short-term directional prediction and showed that textual features can add value when the entity linkage is explicit and the horizon is short. Hagenau, Liebmann, and Neumann then improved on simple bag-of-words pipelines by introducing context-capturing features for financial news classification. Later work replaced manual text features with event embeddings, hierarchical attention, and latent-variable sequence models. Ding et al.\ modeled structured events, Hu et al.\ modeled news sequences with attention, and Xu and Cohen combined tweets with price history in a deep generative framework that reported state-of-the-art performance on a newly collected dataset. The common lesson is that text representation, entity resolution, and temporal alignment are as important as the downstream predictor. \cite{schumaker2009azfin, hagenau2013automated, ding2015event, hu2018chaotic, xu2018stock}

\section{Social Media Sentiment}

Social media has attracted sustained attention because it is timely, large in volume, and rich in investor reaction. Bollen, Mao, and Zeng analyzed approximately 9.8 million tweets from February to December 2008 and linked mood dimensions to subsequent movements in the Dow Jones Industrial Average. Their best reported directional result reached 87.6\% accuracy after including mood information, but the study was conducted at market-index level rather than firm level. That design is important because aggregate mood can correlate with broad market movement without solving the harder task of mapping noisy user text to specific equities. \cite{bollen2011twitter}

Nguyen, Shirai, and Velcin moved closer to firm-level prediction by using company-related message-board data and topic-aware sentiment features. Their experiment covered 18 stocks over one year, with 171 transaction dates for training and 78 for testing. Relative to a price-only baseline, their topic-sentiment method improved mean accuracy by 2.07 percentage points across the 18 stocks and by 9.83 percentage points on the subset of difficult stocks. This is a useful result because it quantifies the value of better ticker-specific filtering. Its limitation is that the setup remains short-horizon, single-language, and dependent on shallow topic extraction rather than contextual language models. \cite{nguyen2015social}

Mao, Counts, and Bollen broadened the comparison by evaluating survey sentiment, news headlines, Twitter data, and search-engine signals in the same predictive study. Their main result was not that every online signal helps, but that different sources behave differently in time. Search and social signals were more informative for short-run market behavior than traditional survey measures. This distinction matters because it suggests that sentiment source is not an interchangeable design choice. Speed, topic granularity, and user noise all affect what a sentiment variable actually measures. \cite{mao2011predicting}

Taken together, the social-media literature shows that sentiment can carry predictive information, but only after heavy filtering. Short text length, sarcasm, reposted content, unofficial company references, and bot-driven amplification all weaken the link between message sentiment and tradable firm-specific information. These constraints are especially relevant for emerging markets with bilingual discussion, because entity linking and language variation become harder before modeling even begins. \cite{bollen2011twitter, nguyen2015social, albahli2022twitter}

Arabic social-media forecasting work is especially relevant to the future direction of this thesis. AlKhatib et al.\ studied Arabic tweets and stock-market changes in the Arab world, including EGX30, using correlation and Granger-causality style analysis. Their study supports the idea that Arabic investor discussion may contain market-relevant information for Egyptian equities. However, it also reinforces the need for careful source access and entity filtering: social-media relevance depends on whether the collected posts can be legally obtained, temporally aligned, and linked to the correct market or company. \cite{alkhatib2016arabictweets}

\section{Financial News Sentiment}

Financial news provides a cleaner signal source than social media because the language is more structured and the company linkage is more explicit. Mohan et al.\ assembled a five-year dataset of daily prices for S\&P 500 firms together with more than 265,000 related financial news articles, then compared deep models built from price data, sentiment information, and hybrid text-price inputs. The scale of the dataset is a strength because many earlier sentiment studies were limited by sparse article counts. The main limitation is external validity: a large U.S.\ large-cap news panel is a much easier information environment than a thinner bilingual EGX news stream. \cite{mohan2019news}

Dahal et al.\ addressed a different problem. Rather than mixing many architectures and feature sets, they compared LSTM and GRU models under identical conditions with and without financial news sentiment. Performance was assessed with RMSE, MAPE, and the correlation coefficient \(R\). Their conclusion was that the sentiment-augmented variants outperformed stock-only variants, which is methodologically useful because it isolates the effect of adding text. At the same time, the study focuses on regression quality for developed-market data, not on leakage-sensitive directional classification for sparse firm-level coverage. \cite{dahal2023comparative}

Ho and Huang fused sentiment analysis with candlestick-chart representations in a two-branch convolutional architecture. Their experiments covered five high-demand stocks, namely Apple, Tesla, IBM, Amazon, and Google, and they evaluated forecast windows of 4, 6, 8, and 10 days. The best reported accuracy was 75.38\% for Apple. This study is useful because it demonstrates that representation fusion can outperform either visual price input or sentiment input alone. Its limitation is equally clear: the result is stock-specific, U.S.-centric, and based on social-media sentiment rather than a firm-linked financial-news archive. \cite{huang2021candlestick}

Cristescu et al.\ used market-news sentiment as an exogenous factor in regression models for stock opening prices. Their most relevant contribution is not a headline accuracy number but the explicit goodness-of-fit analysis. One reported NARX specification achieved an ANOVA significance level of \(p = 0.020\), which supports the claim that sentiment can improve model fit in that setting. The limitation is that regression fit is not the same as robust directional forecasting. A model can improve opening-price estimation while still failing to discriminate reliably between upward and downward moves. \cite{cristescu2022marketnews}

Zhao and Yang proposed a deeper fusion system, SA-DLSTM, that combines an emotion-enhanced CNN, a denoising autoencoder, and an LSTM for stock movement prediction on the Hang Seng Index. The paper is important because it captures several design ideas that recur in later literature: sentiment-strength weighting, feature denoising, and sequence modeling over multivariate financial inputs. Its main limitation is that the target is still an index-level market series. That setting avoids the more difficult problem of article-to-firm relevance matching, which is central in the present thesis. \cite{zhao2023integrated}

Liapis et al.\ provide the most useful benchmark-oriented reference for this thesis. Their study evaluates 30 forecasting methods over 67 feature setups and three forecast horizons. The strongest conclusion is not that one sentiment method wins, but that no single sentiment configuration dominates all horizons. This is an important correction to the literature because it shows that sentiment effectiveness is conditional on representation and horizon choice. It also reinforces the need for conservative reporting. A sentiment setup that performs well for one horizon may be statistically unremarkable or even weak at another. \cite{liapis2023deep}

\begin{table}[ht]
\centering
\scriptsize
\renewcommand{\arraystretch}{1.15}
\setlength{\tabcolsep}{4pt}
\begin{tabular}{>{\raggedright\arraybackslash}p{1.8cm}>{\raggedright\arraybackslash}p{2.7cm}>{\raggedright\arraybackslash}p{1.8cm}>{\raggedright\arraybackslash}p{1.9cm}>{\raggedright\arraybackslash}p{2.1cm}>{\raggedright\arraybackslash}p{2.9cm}}
\toprule
Study & Market and sample & Text input & Model family & Reported result & Main limitation \\
\midrule
Bollen et al.\ \cite{bollen2011twitter} & DJIA, about 9.8 million tweets, Feb.--Nov.\ 2008 & Twitter mood streams & GPOMS mood features with SOFNN & 87.6\% direction accuracy & Index-level setup with no firm-level entity mapping \\
Nguyen et al.\ \cite{nguyen2015social} & 18 stocks, 1 year, 171 train dates and 78 test dates & Message-board topic sentiment & Topic-sentiment features with linear SVM & +2.07 percentage points mean accuracy over price-only; +9.83 on difficult stocks & Short horizon, shallow topic model, single-language data \\
Mohan et al.\ \cite{mohan2019news} & S\&P 500 firms, 5 years, more than 265,000 news articles & Financial news & Deep text, price, and fusion models & Large-scale comparison across price-only, text-only, and fusion settings & Developed-market large-cap panel with no bilingual routing problem \\
Dahal et al.\ \cite{dahal2023comparative} & Developed-market stock data with matched financial news & Financial news plus stock features & LSTM and GRU & Sentiment-augmented models improved RMSE, MAPE, and \(R\) over stock-only variants & Regression emphasis, no firm-level bilingual challenge \\
Ho and Huang \cite{huang2021candlestick} & Five U.S.\ stocks, 4 to 10 day horizons & Social-media sentiment plus candlestick charts & Two-branch CNN & Best reported accuracy 75.38\% for Apple & Stock-specific result and social-media input rather than financial-news input \\
Cristescu et al.\ \cite{cristescu2022marketnews} & Selected stocks with lagged opening-price models & Market news sentiment & Linear and NARX regression & One reported NARX fit significant at \(p = 0.020\) & Goodness of fit does not guarantee robust directional discrimination \\
Zhao and Yang \cite{zhao2023integrated} & Hang Seng Index and public opinion streams & Online sentiment plus market variables & ECNN + DAE + LSTM & Higher accuracy than baselines in the integrated framework & Index-level setting avoids article-to-firm relevance matching \\
Liapis et al.\ \cite{liapis2023deep} & 30 methods, 67 feature setups, 3 horizons & Sentiment time series plus prices & Broad benchmark & No single sentiment setup dominated all horizons & Strong evidence that gains are horizon- and setup-specific \\
\bottomrule
\end{tabular}
\caption{Quantitative comparison of representative sentiment-based stock forecasting studies.}
\label{tab:lit_comparison}
\end{table}

\section{Why Sentiment Does Not Always Improve Forecasting}

The literature does not support a simple claim that sentiment always improves predictive accuracy. Several failure modes recur. The first is semantic coarseness. A positive article can describe a low-impact corporate event, while a seemingly neutral article can contain material financing or governance information. Mudinas, Zhang, and Levene argue that sentiment attitudes alone do not consistently explain future price changes and that any useful relation is strongly case-specific. This is a direct warning against treating polarity as a sufficient event representation. \cite{mudinas2019market}

The second failure mode is temporal misalignment. If article sentiment is joined to prices on the wrong day, or if a target window overlaps the feature window, reported performance can be inflated. Liapis et al.\ show that different sentiment feature setups behave differently across horizons, which implies that loose alignment can easily distort comparisons. In other words, methodological noise can be mistaken for predictive signal. \cite{liapis2023deep}

The third failure mode is data scarcity and imbalance. Sparse firm-level article panels create unstable feature histories, and directional labels often lean toward the majority class. A model can then achieve acceptable raw accuracy by overpredicting the dominant direction. This issue is especially relevant for emerging markets where article volume is much lower than in U.S.\ large-cap datasets. The combination of sparse text and imbalanced labels makes balanced accuracy and confusion-matrix analysis more informative than headline accuracy alone. \cite{shi2021deeplearning, liapis2023deep}

The fourth failure mode is source mismatch. Social media is fast and abundant, but it is noisy and weakly linked to formal firm events. Financial news is cleaner and more structured, but it is slower and less dense. The contrast between Bollen, Nguyen, and Mao on one side and Mohan, Dahal, and Cristescu on the other shows that sentiment source changes both the signal quality and the methodological difficulty. No source is universally superior. The choice must match the research objective, the market, and the entity-linking assumptions. \cite{bollen2011twitter, nguyen2015social, mao2011predicting, mohan2019news, dahal2023comparative, cristescu2022marketnews}

\section{Domain-Specific and Arabic Language Models}

The quality of the sentiment signal depends heavily on the language model used to produce it. Generic sentiment tools often misread financial language because words such as ``liability'', ``dilution'', or ``guidance'' carry domain-specific meanings. FinBERT was designed precisely for this issue. The original paper reports a 14 percentage point improvement over the previous state of the art on Financial PhraseBank, which makes it a much stronger candidate for financial English than a generic sentiment classifier. This result matters because weak sentence-level sentiment estimation will propagate directly into weaker daily forecasting features. \cite{araci2019finbert}

The Arabic side requires an additional specialization step. AraBERT reported state-of-the-art performance on most tested Arabic NLP tasks and the original evaluation reports improvements of up to 2.53 F1 points on the downstream benchmarks. ARBERT and MARBERT extended this line further and introduced the ARLUE benchmark, where their best model achieved an aggregate score of 77.40 and outperformed XLM-R Large despite the latter being substantially larger. These results make a strong case for routing Arabic text through Arabic-specialized encoders rather than treating Arabic as a secondary case inside a generic multilingual pipeline. \cite{antoun2020arabert, abdulmageed2021arbert}

Variant sensitivity is equally important. CAMeLBERT compared Arabic models across five task types and 12 datasets and concluded that variant proximity between pre-training data and downstream data matters more than raw pre-training size. This finding is directly relevant for bilingual EGX forecasting. The problem is not only to classify sentiment, but to classify it in a market where Arabic and English articles occupy different stylistic and domain positions. A single shared model may be simpler to deploy, but the literature suggests that such simplicity often comes at the cost of linguistic fit. \cite{inoue2021camelbert}

At the same time, the Arabic-model literature leaves an open gap. These models are evaluated on Arabic NLP benchmarks, not on bilingual financial forecasting benchmarks for the Egyptian market. Their strong benchmark performance justifies their use as sentiment encoders, but it does not by itself prove that they will improve stock prediction. That remaining empirical question is one of the contributions of the present thesis. \cite{antoun2020arabert, abdulmageed2021arbert, inoue2021camelbert}

Recent Arabic financial NLP resources make this gap more visible. SAHM introduces Arabic financial and Shari'ah-compliant reasoning tasks, while recent Saudi-market work studies Arabic financial sentiment with large language models. These papers support the need for Arabic finance-specific evaluation rather than relying only on general Arabic sentiment benchmarks. They also clarify the limitation of the present thesis: without independent human-labeled Arabic financial sentiment examples, AraBERT can be adapted to financial text through masked-language modeling, but it cannot be validly evaluated as a supervised Arabic FinBERT equivalent. \cite{elbadry2026sahm, albaqawi2026arabicfinancialsentiment}

\section{Evaluation Metrics and Research Gap}

Comparison across prior forecasting studies is complicated by inconsistent evaluation practice. Regression-oriented papers often report RMSE, MAE, MAPE, or correlation. Directional papers more often report accuracy, precision, recall, and F1. A few papers add simulated returns, which is useful but depends strongly on trading assumptions. These metrics are not interchangeable. A classifier can show a favorable raw accuracy while failing to recover the minority class. A regression model can reduce RMSE while still predicting the sign of the move poorly. For this reason, metrics must be interpreted in relation to the target definition and the class balance. \cite{dahal2023comparative, liapis2023deep, shi2021deeplearning, zhao2023integrated}

The gap addressed by this thesis is therefore both geographical and methodological. Most sentiment-based forecasting studies focus on English text, developed markets, single-language pipelines, or index-level prediction. Very few studies examine bilingual Arabic-English financial-news routing for Egyptian equities, and fewer still evaluate that problem under conservative walk-forward validation with explicit leakage and overfitting checks. The present work narrows the scope to three EGX tickers and a single bilingual financial-news source so that article relevance, language routing, and temporal alignment can be studied under cleaner conditions. The contribution is not to assume that sentiment is universally powerful. The contribution is to test, under a more locally relevant and more tightly controlled design, how much forecasting value bilingual financial-news sentiment actually carries.

\section{Chapter Summary}

This chapter reviewed the state of the art in sentiment-based stock forecasting and identified the main lessons that shape the present thesis. Text can help, but the reported gains depend on source quality, feature design, horizon choice, and validation rigor. Financial-news studies are usually cleaner than social-media studies, yet they remain concentrated in English and in developed markets. Domain-specific financial English models and Arabic-specialized transformers offer strong upstream sentiment encoders, but they have rarely been evaluated together in a bilingual EGX forecasting setting. These gaps motivate the methodological choices developed in the next chapter.
